% !TEX root = ../main.tex
\section{Marco teórico y ejercicios resueltos}

\subsection{Comparativa formal de complejidades teóricas}
La Tabla \ref{tab:complejidad_teorica} resume las cotas asintóticas para las operaciones fundamentales entre un montículo binario estándar y un montículo de Fibonacci.

\begin{table}[H]
    \caption{Comparativa de complejidades temporales entre estructuras de prioridad}
    \label{tab:complejidad_teorica}
    \centering
    \small
    \begin{tabular*}{\textwidth}{@{\extracolsep{\fill}}lcc}
        \toprule
        \textbf{Operación} & \textbf{Montículo Binario (\texttt{heapq})} & \textbf{Montículo de Fibonacci} \\
        \midrule
        \texttt{insert} & $O(\log n)$ (peor caso) & $O(1)$ (amortizado) \\
        \texttt{minimum} & $O(1)$ (peor caso) & $O(1)$ (peor caso) \\
        \texttt{extract-min} & $O(\log n)$ (peor caso) & $O(\log n)$ (amortizado) \\
        \texttt{decrease-key} & $O(\log n)$ (con búsqueda) & $O(1)$ (amortizado) \\
        \texttt{union} & $O(n)$ o no constante & $O(1)$ (peor caso) \\
        \texttt{delete} & $O(\log n)$ & $O(\log n)$ (amortizado) \\
        \bottomrule
    \end{tabular*}
\end{table}

% -------------------------------------------------------------------------
\subsection{Ejercicio resuelto 1: instrumentación de la función potencial}
Implementa la contabilidad del potencial y el cálculo del costo amortizado formal.

\begin{lstlisting}[language=Python, caption={Código del Ejercicio Resuelto 1 (función potencial y costo amortizado)}, label={lst:resuelto_1}]
from dataclasses import dataclass

@dataclass(frozen=True)
class EstadoPotencial:
    arboles: int
    marcados: int
    def potencial(self) -> int: return self.arboles + 2 * self.marcados

def costo_amortizado(costo_real, antes: EstadoPotencial, despues: EstadoPotencial):
    delta = despues.potencial() - antes.potencial()
    return costo_real + delta, delta

# Caso de prueba oficial
antes = EstadoPotencial(arboles=4, marcados=3)
despues = EstadoPotencial(arboles=7, marcados=0)
amortizado, delta = costo_amortizado(4, antes, despues)
print(f"Phi antes: {antes.potencial()} | Phi despues: {despues.potencial()} | Delta: {delta} | Amortizado: {amortizado}")
assert antes.potencial() == 10 and despues.potencial() == 7 and delta == -3 and amortizado == 1
\end{lstlisting}

\begin{lstlisting}[style=consolestyle, caption={Salida de consola del Ejercicio Resuelto 1}]
Phi antes: 10 | Phi despues: 7 | Delta: -3 | Amortizado: 1
\end{lstlisting}

El costo real fue $c_i = 4$, pero la liberación de potencial $\Delta \Phi = -3$ resultó en un costo amortizado contable de $\hat{c}_i = c_i + \Delta \Phi = 4 - 3 = 1$.

% -------------------------------------------------------------------------
\subsection{Ejercicio resuelto 2: consolidación programada por grados}
Modela la fusión de raíces por grados iguales mediante una tabla de colisiones.

\begin{lstlisting}[language=Python, caption={Código del Ejercicio Resuelto 2 (consolidación simplificada)}, label={lst:resuelto_2}]
from dataclasses import dataclass, field

@dataclass
class Raiz:
    clave: int; grado: int = 0; hijos: list = field(default_factory=list)

def enlazar(a, b):
    padre, hijo = (a, b) if a.clave <= b.clave else (b, a)
    padre.hijos.append(hijo); padre.grado += 1
    return padre

def consolidar(raices):
    por_grado = {}
    for raiz in raices:
        actual = raiz
        while actual.grado in por_grado:
            otra = por_grado.pop(actual.grado)
            actual = enlazar(actual, otra)
        por_grado[actual.grado] = actual
    return list(por_grado.values())

raices = [Raiz(23), Raiz(7), Raiz(21), Raiz(18, 1), Raiz(52), Raiz(38, 1), Raiz(17, 1), Raiz(24, 2)]
resultado = consolidar(raices)
print([(r.clave, r.grado) for r in resultado])
assert len(resultado) == len(set(r.grado for r in resultado)) and min(r.clave for r in resultado) == 7
\end{lstlisting}

\begin{lstlisting}[style=consolestyle, caption={Salida de consola del Ejercicio Resuelto 2}]
[(52, 0), (23, 1), (17, 2), (7, 3)]
\end{lstlisting}

Al concluir, los árboles resultantes poseen grados disjuntos ($0, 1, 2, 3$), y el mínimo global ($7$) se preserva en la raíz de grado mayor.

% -------------------------------------------------------------------------
\subsection{Ejercicio resuelto 3: simulación de decrease-key con heapq}
Implementa una cola de prioridad basada en \texttt{heapq} con invalidación perezosa.

\begin{lstlisting}[language=Python, caption={Código del Ejercicio Resuelto 3 (cola de prioridad con heapq perezoso)}, label={lst:resuelto_3}]
import heapq, itertools
REMOVED = object()

class HeapqPriorityQueue:
    def __init__(self):
        self.heap, self.entries, self.counter = [], {}, itertools.count()

    def insert(self, item, priority):
        if item in self.entries: self.remove(item)
        entry = [priority, next(self.counter), item]
        self.entries[item] = entry
        heapq.heappush(self.heap, entry)

    def remove(self, item):
        entry = self.entries.pop(item)
        entry[2] = REMOVED

    def decrease_key(self, item, new_priority):
        if new_priority > self.entries[item][0]: raise ValueError("Clave mayor")
        self.insert(item, new_priority)

    def extract_min(self):
        while self.heap:
            priority, _, item = heapq.heappop(self.heap)
            if item is not REMOVED:
                del self.entries[item]
                return item, priority
        raise IndexError("Cola vacia")

pq = HeapqPriorityQueue()
pq.insert("a", 20); pq.insert("b", 12); pq.insert("c", 30)
pq.decrease_key("c", 5)
print(f"Minimo: {pq.extract_min()} | Activos: {len(pq.entries)} | Tamano fisico heap: {len(pq.heap)}")
\end{lstlisting}

\begin{lstlisting}[style=consolestyle, caption={Salida de consola del Ejercicio Resuelto 3}]
Minimo: ('c', 5) | Activos: 2 | Tamano fisico heap: 3
\end{lstlisting}

La diferencia entre el tamaño físico (3) y los elementos activos (2) evidencia las entradas obsoletas que permanecen en memoria hasta ser retiradas en el tope.
